% Generated draft: requires mathematical and source review.
\documentclass[11pt]{article}
\usepackage[margin=25mm]{geometry}
\usepackage{fontspec}
\setmainfont{Latin Modern Roman}
\usepackage{amsmath,amssymb,mathtools,unicode-math}
\setmathfont{Latin Modern Math}
\usepackage{enumitem,needspace,xurl,hyperref}
\setlength{\parindent}{0pt}
\setlength{\parskip}{5pt}
\setlength{\emergencystretch}{3em}
\newcommand{\N}{\mathbb{N}}
\newcommand{\Z}{\mathbb{Z}}
\newcommand{\Q}{\mathbb{Q}}
\newcommand{\R}{\mathbb{R}}
\newcommand{\C}{\mathbb{C}}
\newcommand{\F}{\mathbb{F}}
\newcommand{\A}{\mathbb{A}}
\newcommand{\PP}{\mathbb{P}}
\newcommand{\E}{\mathbb{E}}
\newcommand{\eps}{\varepsilon}
\newcommand{\dd}{\,\mathrm{d}}
\begin{document}
\section*{UnsolvedMath ID: 2800103}
\textbf{Source problem code:} catalogue entry
\subsection{Definitions and mathematical statement}
For each $n\geq1$, let $W^{(n)}$ be a real symmetric Wigner matrix: the upper-triangular entries $W_{ij}$, $1\leq i\leq j\leq n$, are independent $\mathcal N(0,1)$ variables and $W_{ji}=W_{ij}$. This makes precise the usual independent-entry interpretation of the source's symmetric Gaussian matrix; the diagonal variance is kept at the stated value $1$.

For a real symmetric matrix $B\in\R^{n\times n}$, define
\[
 Q_n(B)=\max\{\operatorname{tr}(BX):X\in\R^{n\times n}\text{ symmetric},\ X\succeq0,\ X_{ii}=1\ (1\leq i\leq n)\}.
\]
There is no rank constraint on $X$. With $\mathbf1_n$ the all-ones column vector and $\xi\geq0$ fixed as $n\to\infty$, the source defines
\[
 q(\xi)=\lim_{n\to\infty}\frac1n\E\,Q_n\left(
 \frac\xi n\mathbf1_n\mathbf1_n^{\mathsf T}+\frac1{\sqrt n}W^{(n)}\right),
\]
where expectation is over $W^{(n)}$. The source uses the limit as a definition and supplies no existence proof in this excerpt.

\textbf{Open Problem 1.3 (Afonso S. Bandeira).} Determine exactly
\[
 \xi_* =\inf\{\xi\geq0:q(\xi)>2\}.
\]
The comparison is strict, the baseline is exactly $2$, and the usual extended-real convention gives $\inf\varnothing=+\infty$.

The supplied research summary says that, as of its 2026-08-05 search, the exact threshold was not confirmed resolved and no verifiable sharp-theorem or exact-value citation was located. It describes a heuristic transition at a signal of constant order one, by comparison with BBP/spiked-Wigner semidefinite-program phase transitions. This heuristic does not specify a value and is not substituted for a theorem.
\subsection{Short English statement}
Find the exact signal strength at which the limiting expected semidefinite optimum for a rank-one-spiked Gaussian Wigner matrix rises strictly above $2$ per dimension. The source gives the limiting formula and an unverified heuristic, without an exact-threshold theorem.
\subsection{Sources}
\begin{itemize}
\item Local UnsolvedMath registry, permanent ID 2800103, source code catalogue entry, ``10 Lectures and 42 Open Problems --- Open Problem 1.3''; proposed by Afonso S. Bandeira in 2015. Supplied status open, classification OPEN-TRIAGE, updated 2026-08-06; research-summary search date 2026-08-05.
\item Bandeira, source list ``42 Open Problems in Data Science'' (2016), Open Problem 1.3; supplied author-blog URL \url{https://afonsobandeira.wordpress.com/2015/09/18/10l42ppca/}. The record reports HTML extraction and access on 2026-07-29. It says the source presents an open item and current status requires release review. Rights: NEEDS\_REVIEW; publicly accessible author blog.
\item The research summary mentions BBP and spiked-Wigner SDP comparisons without a precise supporting citation, and expressly reports no verified exact-value theorem found.
\end{itemize}
\end{document}
